% additional packages
\usepackage{xspace}
% \usepackage{graphicx}  % Already loaded in main.tex
% \usepackage{subfigure}
% \usepackage{amsmath}  % Already loaded in main.tex
% \usepackage{amssymb}  % Already loaded in main.tex
% \usepackage{amsthm}   % Already loaded in main.tex
\usepackage{adjustbox}
% \usepackage{algorithm}
% \usepackage{algpseudocode}
\usepackage{wrapfig}
% \usepackage{caption}
\usepackage{lipsum}
% \usepackage{booktabs}  % Already loaded in main.tex
\usepackage{multirow}
\usepackage{makecell}
\usepackage{pifont}

% ---------- Markers ----------
\newcommand{\cmark}{\ding{51}}% ✓
\newcommand{\xmark}{\ding{55}}% ✗

% ---------- PREAMBLE ----------
% Note: algorithm and algorithmic packages are already loaded by icml2026.sty
% \usepackage{algorithm}
% \usepackage[noend]{algpseudocode}  % algorithmicx + pseudocode style
% \usepackage{algpseudocode}  % Commented out - conflicts with algorithmic from icml2026.sty

% Rename the default keywords (using algorithmic package syntax, not algpseudocode)
\renewcommand{\algorithmicrequire}{\textbf{Input:}}
\renewcommand{\algorithmicensure}{\textbf{Output:}}
% --------------------------------

% \usepackage{subcaption}  % Already loaded in main.tex

\newcommand{\name}[0]{TaH\xspace}
\newcommand{\xhdr}[1]{{\noindent\bfseries #1}.}
\newcommand{\xhdrq}[1]{{\noindent\bfseries #1}}
\newcommand{\xhdrl}[1]{{\noindent\bfseries\underline{#1}}.}
\newcommand{\rom}[1]{\textup{\uppercase\expandafter{\romannumeral#1}}}

\newcommand{\highlightblock}[1]{\begin{center}\vspace{-0.1cm}\emph{#1}\vspace{-0.1cm}\end{center}}

\newcounter{observation}
\newenvironment{observation}[1][]{\refstepcounter{observation}\par\medskip
   \noindent \textbf{Observation~\theobservation. #1} \rmfamily}{\medskip}

\newcounter{design}
\newenvironment{design}[1][]{\refstepcounter{design}\par\medskip
   \noindent \textbf{Design Choice~\thedesign. #1} \rmfamily}{\medskip}

%%%%%%%%%%%%%%%%%%%% begin color box example %%%%%%%%%%%%%%%%%%%%
% package for example
\usepackage[most]{tcolorbox}
\usepackage{colortbl}
\usepackage{pifont}
\usepackage{makecell}

\tcbuselibrary{skins, breakable}   % already loaded if you used it before
\usepackage{tabularx, booktabs}    % flexible table + nice rules

%%%%%%% color %%%%%%%
\definecolor{darkred}{rgb}{0.5,0,0}
\definecolor{darkgreen}{rgb}{0,0.5,0}
\definecolor{tahgray}{RGB}{242,242,242}

%%%%%%% color box %%%%%%%

%%% green color box %%%
\newtcolorbox[
  auto counter,
  list inside=examplelist,
  % number within=section
]{greenbox}[2][]{
  enhanced,
  title={Example~\thetcbcounter. \textbf{#1}},
  colback=green!5!white,
  colframe=green!50!black,
  colbacktitle=green!60!black,
  coltitle=white,
  #2
}

%%% blue color box %%%
\newtcolorbox[
  auto counter,
  list inside=examplelist
]{bluebox}[2][]{
  enhanced,
  title={Text~\thetcbcounter. \textbf{#1}},
  colback=cyan!5!white,  % Light blue background
  colframe=cyan!50!black,  % Darker blue frame
  colbacktitle=cyan!75!black,  % Even darker blue title background
  coltitle=white,
  #2
}

% --- handy 7‑argument wrapper: #1–#7 as shown below ---
\newcommand{\threecompare}[7]{%
  \begin{tcolorbox}[
      enhanced,
      sharp corners,
      fonttitle=\bfseries\large,
      colback=cyan!6!white,       % body background
      colframe=cyan!50!black,     % frame
      colbacktitle=cyan!80!black, % title background
      coltitle=white,             % title text
      title={#1},                 % #1 = question
      drop shadow               % subtle shadow
    ]
    \setlength{\extrarowheight}{2pt}      % vertical breathing room
    \renewcommand{\arraystretch}{1.15}

    % ---- 3‑column table with dashed vertical separators ----
    \begin{tabularx}{\linewidth}{>{\raggedright\arraybackslash}X
      : >{\raggedright\arraybackslash}X
      : >{\raggedright\arraybackslash}X}
      \rowcolor{cyan!15}          % header band
      \textbf{#2} & \textbf{#3} & \textbf{#4} \\[4pt]
      #5 & #6 & #7
    \end{tabularx}
  \end{tcolorbox}
}
% ------------------------------------------------------
% Usage:
% \threecolbluebox{<question>}{<title1>}{<title2>}{<title3>}
%                    {<answer1>}{<answer2>}{<answer3>}
% ------------------------------------------------------


%%%%%%% code %%%%%%%

%%% python code %%%
\lstset{
  language=Python,
  basicstyle=\ttfamily\scriptsize,
  breaklines=true,
  tabsize=4,
  numbers=left,
  moredelim=[is][\color{blue}]{|}{|}
}

%%%%%%% usage templates %%%%%%%

%%% using in-place text version %%%
\newcommand{\imgThreeTextConv}[9]{
\begin{bluebox}[#1]{#9}

\begin{center}
\includegraphics[width=\textwidth,height=10cm,keepaspectratio]{#3} % Question figure
\end{center}

\textbf{Question:}

#2 % Question text

\tcbline
\textbf{Original Response:}

#4 % FP16 response

\tcbline
\textbf{#5 Response:}

#6 % Second response

\tcbline
\textbf{#7 Response:}

#8 % Third response

\end{bluebox}
}


%%% using latex version %%%
\newcommand{\imgThreeLatexConv}[9]{
\begin{bluebox}[#1]{#9}
\textbf{Question:}
% \scriptsize
\input{#2} % Question text

\begin{center}
\includegraphics[width=\textwidth,height=10cm,keepaspectratio]{#3} % Question figure
\end{center}

\tcbline
\textbf{FP16 Response:}
% \scriptsize
\input{#4} % FP16 response

\tcbline
\textbf{#5 Response:}
% \scriptsize
\input{#6} % AWQ W3 response

\tcbline
\textbf{#7 Response:}
% \scriptsize
\input{#8} % AWQ W3 response


\end{bluebox}
}

%%%%%%%%%%%%%%%%%%%% end color box example %%%%%%%%%%%%%%%%%%%%

% Theorem definitions moved to main.tex to avoid conflicts
% \newtheorem{theorem}{Theorem}
% \newtheorem{corollary}{Corollary}[theorem] % Use theorem counter as `parent`
% \newtheorem{definition}{Definition}[section]
% \newtheorem{lemma}{Lemma}[section]

\newcommand{\tianyu}[1]{{\color{blue}[\textbf{tianyu}: #1]}}
\newcommand{\xuefei}[1]{{\color{cyan}[\textbf{xuefei}: #1]}}
\newcommand{\yichen}[1]{\textcolor{orange}{[\textbf{yichen}: #1]}}
\newcommand{\zekai}[1]{{\color{teal}[\textbf{zekai}: #1]}}

\newcommand{\RA}[1]{{\color{teal}#1}}
\newcommand{\RB}[1]{{\color{orange}#1}}
\newcommand{\RC}[1]{{\color{cyan}#1}}
\newcommand{\CQ}[1]{{\color{brown}#1}}

% \newcommand{\todo}[1]{{\color{magenta}[\textbf{TODO}: #1]}}  % Commented out - using todonotes package instead
\newcommand{\fix}{\marginpar{FIX}}
\newcommand{\new}{\marginpar{NEW}}

\newcommand{\revision}[1]{\textcolor{blue}{#1}}
\newcommand{\typo}[1]{\textcolor{orange}{#1}}

% Efficiency analysis shorthands
\newcommand{\F}{\mathrm{FLOPs}}
\newcommand{\M}{\mathrm{Mem}}




% --- disable draft comments for submission ---
% \renewcommand{\todo}[1]{}
% \renewcommand{\tianyu}[1]{}
% \renewcommand{\yichen}[1]{}
% \renewcommand{\zekai}[1]{}
% \renewcommand{\xuefei}[1]{}